\rtmtight
\begin{tblr}{width=\textwidth,colspec={|Q[c,m,2.1cm]|X[l,m]|},hlines,vlines,hspan=minimal,rowsep=1.5pt,colsep=3pt,row{1}={bg=headgray,font=\bfseries}}
\SetCell{c,m}\hfil\logo\hfil & \textbf{POLITEKNIK NEGERI MALANG}\par\textbf{JURUSAN TEKNOLOGI INFORMASI}\par\textbf{PROGRAM STUDI: D4 TEKNIK INFORMATIKA} \\
\SetCell[c=2]{l} \textbf{RENCANA TUGAS MAHASISWA (RTM) DAN RUBRIK PENILAIAN} & \\
Nama Mata Kuliah & Pembelajaran Mesin \\
Kode Mata Kuliah & RTI255004 \quad \textbf{sks / Jam perminggu:} 3 SKS/6 jam \quad \textbf{Semester:} 5 \\
Dosen Pengampu & Tim Pengajar Program Studi D4 TI \\
\end{tblr}
\assessmentblock{Tugas Implementasi Algoritma ML}{Case Method dan praktik ML}{SCPMK0706-04001, SCPMK1005-04002}{Mahasiswa mengimplementasikan dan mengevaluasi algoritma ML untuk menyelesaikan studi kasus klasifikasi, regresi, atau clustering menggunakan Python/scikit-learn.}{Menganalisis dataset, memilih algoritma yang tepat, mengimplementasikan pipeline ML, mengevaluasi kinerja, dan mendokumentasikan keputusan.}{Notebook Python, laporan implementasi, hasil evaluasi model, dan/atau bahan presentasi.}{Ketepatan implementasi algoritma ML & 12 & Algoritma diimplementasikan dengan benar, pipeline lengkap (preprocessing, training, prediction), dan output dapat direproduksi. & Implementasi sebagian besar benar, tetapi beberapa langkah pipeline atau parameterisasi belum optimal. & Implementasi salah atau pipeline tidak lengkap sehingga hasil tidak valid. \\ \\
Analisis dan perbandingan model & 10 & Minimal dua algoritma dibandingkan menggunakan metrik yang tepat, kelebihan dan kekurangan dianalisis secara kritis. & Perbandingan dilakukan, tetapi analisis metrik atau argumentasi pilihan model masih kurang mendalam. & Tidak ada perbandingan yang bermakna atau metrik tidak tepat untuk jenis masalah. \\ \\
Dokumentasi kode dan laporan & 6 & Notebook terdokumentasi lengkap dengan penjelasan setiap langkah, kode bersih, dan laporan menyimpulkan temuan dengan jelas. & Dokumentasi tersedia, tetapi penjelasan langkah atau kesimpulan masih kurang lengkap. & Kode tidak terdokumentasi atau laporan tidak menjelaskan proses dan temuan analisis. \\}

\assessmentblock{Kuis ML}{Kuis}{SCPMK0706-04001, SCPMK1005-04002}{Kuis tertulis untuk mengukur pemahaman konsep supervised learning, evaluasi model, dan teknik optimasi ML.}{Menjawab soal pilihan ganda dan/atau uraian terkait algoritma ML, metrik evaluasi, cross-validation, dan hyperparameter tuning.}{Jawaban tertulis kuis.}{Penguasaan konsep supervised dan unsupervised learning & 6 & Konsep dan perbedaan supervised/unsupervised learning, serta cara kerja algoritma utama dijelaskan dengan benar. & Sebagian besar konsep dipahami, tetapi penjelasan cara kerja beberapa algoritma masih kurang akurat. & Konsep tidak dipahami atau cara kerja algoritma dijelaskan secara keliru. \\ \\
Analisis metrik evaluasi dan optimasi model & 4 & Metrik (accuracy, precision, recall, F1, MSE) dipilih tepat untuk jenis masalah dan interpretasi cross-validation serta regularisasi benar. & Metrik dan optimasi sebagian besar tepat, tetapi interpretasi atau pemilihan metrik untuk kasus tertentu masih keliru. & Metrik tidak tepat untuk jenis masalah atau konsep regularisasi tidak dipahami. \\}

\assessmentblock{UTS Supervised Learning dan Evaluasi}{UTS berbasis praktik ML}{SCPMK0706-04001, SCPMK1005-04002}{Evaluasi tengah semester untuk mengukur kemampuan implementasi dan evaluasi algoritma supervised learning menggunakan Python.}{Mengimplementasikan algoritma supervised learning, mengevaluasi model dengan metrik yang tepat, dan membandingkan beberapa pendekatan.}{Notebook Python, output evaluasi model, dan laporan singkat.}{Implementasi supervised learning & 7 & Algoritma supervised learning diimplementasikan dengan benar, data di-split dengan tepat, dan model dilatih tanpa data leakage. & Implementasi sebagian besar benar, tetapi data splitting atau potensi leakage masih perlu perbaikan. & Algoritma tidak diimplementasikan dengan benar atau terdapat data leakage yang signifikan. \\ \\
Evaluasi dan komparasi model & 5 & Model dievaluasi dengan metrik yang tepat, cross-validation diterapkan, dan hasil dibandingkan secara kuantitatif. & Evaluasi dilakukan tetapi pemilihan metrik atau penerapan cross-validation masih kurang tepat. & Evaluasi tidak sistematis atau metrik tidak sesuai dengan jenis masalah. \\ \\
Kualitas analisis dan kesimpulan & 3 & Analisis menjelaskan alasan pemilihan model terbaik dan kesimpulan didukung data evaluasi yang kuat. & Analisis tersedia, tetapi argumentasi pemilihan model atau kesimpulan masih kurang kuat. & Tidak ada analisis yang menjelaskan keputusan pemilihan model. \\}

\assessmentblock{PBL End-to-End ML Project}{Project Based Learning}{SCPMK0706-04001, SCPMK1005-04002}{Mahasiswa membangun solusi ML end-to-end untuk masalah nyata mencakup EDA, feature engineering, pemodelan, optimasi, dan interpretasi.}{Memilih dataset, melakukan EDA, feature engineering, melatih dan mengoptimalkan model, menginterpretasikan hasil, dan mempresentasikan proyek.}{Notebook Python lengkap, laporan proyek ML, visualisasi, dan bahan presentasi.}{Kualitas pipeline ML end-to-end & 10 & Pipeline mencakup EDA, preprocessing, feature engineering, modeling, dan evaluasi secara sistematis dan dapat direproduksi. & Pipeline mencakup tahap utama, tetapi beberapa langkah seperti feature engineering atau reproduktibilitas masih kurang. & Pipeline tidak lengkap atau tidak sistematis sehingga hasil tidak dapat diandalkan. \\ \\
Kualitas model dan optimasi & 9 & Hyperparameter tuning dilakukan secara terstruktur, metrik final membuktikan peningkatan, dan regularisasi diterapkan dengan tepat. & Optimasi dilakukan, tetapi metode tuning atau dampaknya pada metrik belum terdokumentasi dengan baik. & Tidak ada optimasi yang bermakna atau performa model tidak dapat dibandingkan sebelum dan sesudah tuning. \\ \\
Presentasi dan dokumentasi proyek & 6 & Proyek dipresentasikan secara komprehensif, insight ML dikomunikasikan dengan jelas, dan pertanyaan teknis dijawab dengan tepat. & Presentasi cukup jelas, tetapi beberapa insight atau jawaban teknis masih kurang mendalam. & Presentasi tidak mencerminkan pemahaman proyek atau tidak dapat menjawab pertanyaan teknis dasar. \\}

\assessmentblock{UAS Evaluasi dan Presentasi ML Project}{UAS berbasis presentasi dan defense}{SCPMK1005-04002}{Mahasiswa mempresentasikan, mengevaluasi, dan mempertahankan hasil proyek ML yang telah dioptimalkan.}{Menyiapkan presentasi final proyek ML, mendemonstrasikan performa model, menginterpretasikan hasil, dan mempertahankan keputusan metodologis.}{Notebook Python final, slide presentasi, laporan evaluasi model, dan visualisasi SHAP/LIME.}{Ketepatan evaluasi dan interpretasi model & 10 & Metrik evaluasi dipilih tepat, interpretabilitas model (SHAP/LIME) diterapkan, dan temuan diinterpretasikan secara akurat dan kontekstual. & Evaluasi dilakukan dengan benar, tetapi interpretabilitas atau kontekstualisasi temuan masih kurang mendalam. & Evaluasi tidak tepat atau tidak ada upaya interpretasi model yang bermakna. \\ \\
Kualitas dan inovasi solusi ML & 7 & Solusi menunjukkan pemilihan model yang tepat, optimasi yang terukur, dan pendekatan yang kreatif dalam menyelesaikan masalah. & Solusi fungsional dan teroptimasi, tetapi inovasi atau justifikasi pendekatan masih kurang kuat. & Solusi tidak dioptimalkan atau pendekatan tidak sesuai dengan karakteristik masalah. \\ \\
Presentasi dan pertahanan argumen & 5 & Presentasi terstruktur, visualisasi mendukung narasi, dan mahasiswa mampu mempertahankan setiap keputusan metodologis dengan meyakinkan. & Presentasi cukup baik, tetapi beberapa keputusan metodologis belum dapat dipertahankan dengan argumen yang kuat. & Mahasiswa tidak mampu menjelaskan keputusan teknis atau presentasi tidak mencerminkan pemahaman proyek. \\}

\begin{tblr}{width=\textwidth,colspec={|Q[l,m,3.2cm]|X[l,m]|},hlines,vlines,hspan=minimal,row{1-3}={bg=headgray},rowsep=1.5pt,colsep=3pt}
Jadwal Pelaksanaan: & Minggu 1--17 sesuai RPS. Setiap evaluasi memiliki RTM dan rubrik multi-kriteria. \\
Lain-Lain Yang Diperlukan: & LMS, repositori/notebook Python, perangkat lunak sesuai mata kuliah, dan perangkat presentasi. \\
Pustaka: & Mengacu pustaka utama dan pendukung pada bagian RPS. \\
\end{tblr}
